%=============================================================================
% Chapter 5: Technology Stack
%=============================================================================
\chapter{Technology Stack}

This chapter documents every technology used in FloodEye, explaining the purpose of each component and the rationale for its selection.

\section{Frontend Technologies}

\subsection{Next.js 16 (App Router)}
\textbf{Purpose:} Full-stack React framework providing server-side rendering, API routes, file-system-based routing, and seamless Vercel deployment.\\
\textbf{Why chosen:} Next.js App Router colocates server components, client components, and API handlers in a single project. Its ISR (Incremental Static Regeneration) with the \texttt{revalidate} export allows fine-grained cache control per API route without a separate caching layer. The Vercel deployment integration is zero-configuration.\\
\textbf{Version:} 16.2.10

\subsection{React 19}
\textbf{Purpose:} UI component library; the foundation of all dashboard views.\\
\textbf{Why chosen:} React 19 introduces concurrent features and improved suspense semantics. The Context API (\texttt{TelemetryProvider}) serves as lightweight global state management without requiring Redux or Zustand.

\subsection{TypeScript 5}
\textbf{Purpose:} Static typing for all source files (\texttt{.\allowbreak{}ts}, \texttt{.\allowbreak{}tsx}).\\
\textbf{Why chosen:} TypeScript eliminates entire categories of runtime errors in the type-rich data pipeline (ThingSpeak field parsing, TelemetryData objects, alert severity enums). The \texttt{FloodPredictionResult}, \texttt{TelemetryData}, \texttt{DashboardSettings}, and \texttt{Alert} interfaces are all strongly typed.

\subsection{Tailwind CSS v4}
\textbf{Purpose:} Utility-first CSS framework for all component styling.\\
\textbf{Why chosen:} Tailwind v4's JIT (just-in-time) compiler generates minimal CSS bundles. The glassmorphism design language (\texttt{bg-white/\allowbreak{}40 backdrop-blur-xl border-white/\allowbreak{}60}) is cleanly expressed with utility classes without custom CSS files.

\subsection{shadcn/ui}
\textbf{Purpose:} Accessible, unstyled component primitives for buttons, dialogs, and form elements.\\
\textbf{Why chosen:} shadcn/ui components are installed directly into the project source rather than as an opaque dependency, making them fully customisable.

\subsection{Recharts 3}
\textbf{Purpose:} SVG-based charting library for all time-series line charts.\\
\textbf{Why chosen:} Recharts integrates naturally with React's rendering model. It provides \texttt{ResponsiveContainer}, \texttt{LineChart}, \texttt{Tooltip}, and \texttt{CartesianGrid} out-of-the-box, which are sufficient for FloodEye's chart requirements without the configuration overhead of D3.

\subsection{MapLibre GL JS 5 + MapTiler}
\textbf{Purpose:} WebGL-accelerated interactive map rendering.\\
\textbf{Why chosen:} MapLibre GL is an open-source fork of Mapbox GL with no usage-based pricing caps. MapTiler provides high-quality raster and vector tile services compatible with MapLibre's style spec. The \texttt{react-map-gl} wrapper provides React-native integration.

\subsection{Framer Motion 12}
\textbf{Purpose:} Animation library for UI micro-animations, card hover effects, and page transitions.\\
\textbf{Why chosen:} Framer Motion provides a declarative API (\texttt{motion.\allowbreak{}div}, \texttt{AnimatePresence}) that integrates directly with React's component tree.

\subsection{GSAP 3 + @gsap/react}
\textbf{Purpose:} ScrollTrigger-based animations on the landing page.\\
\textbf{Why chosen:} GSAP is the industry standard for high-performance scroll-driven animations. The \texttt{ScrollytellingSections} component uses GSAP's \texttt{fromTo} and \texttt{ScrollTrigger} to choreograph section entrance animations.

\subsection{Lenis 1.3}
\textbf{Purpose:} Smooth, physics-based scroll inertia on the landing page.\\
\textbf{Why chosen:} Lenis provides buttery-smooth native scroll override that pairs naturally with GSAP ScrollTrigger.

\subsection{Lucide React}
\textbf{Purpose:} Icon library providing \texttt{Thermometer}, \texttt{Droplets}, \texttt{Wind}, \texttt{Mountain}, \texttt{Ruler}, \texttt{WifiOff}, \texttt{AlertTriangle}, and other icons used throughout the dashboard.

\subsection{jsPDF + jspdf-autotable}
\textbf{Purpose:} Client-side PDF generation for the sensor history export on the History page.\\
\textbf{Why chosen:} No server is required; the PDF is generated entirely in the browser and downloaded directly.

\section{Machine Learning Technologies}

\subsection{TensorFlow 2 / Keras (Training)}
\textbf{Purpose:} Python deep learning framework used to define, train, and export the flood prediction model.\\
\textbf{Role:} The model is defined in \texttt{Model/\allowbreak{}nb\_\allowbreak{}code.\allowbreak{}py} (extracted from \texttt{FloodMonitoring.\allowbreak{}ipynb}). Training runs in Google Colab with GPU acceleration.

\subsection{TensorFlow.js 4.22 (Inference)}
\textbf{Purpose:} JavaScript port of TensorFlow; runs the trained model in the browser using WebGL acceleration.\\
\textbf{Why chosen:} Client-side inference eliminates the need for a dedicated ML serving infrastructure. The model runs on the user's device with no server cost for inference.\\
\textbf{Integration:} The custom \texttt{AttentionLayer} is reimplemented in TypeScript and registered with \texttt{tf.\allowbreak{}serialization.\allowbreak{}registerClass()} before model loading.

\subsection{scikit-learn StandardScaler}
\textbf{Purpose:} Feature normalisation during training.\\
\textbf{Role:} Scaler parameters (mean and scale vectors) are exported to \texttt{scalers.\allowbreak{}json} and loaded by the browser at inference time to apply the same normalisation.

\subsection{TensorFlow.js Converter}
\textbf{Purpose:} Converts the trained Keras \texttt{.\allowbreak{}h5} model to TensorFlow.js LayersModel format (\texttt{model.\allowbreak{}json} + \texttt{group1-shard1of1.\allowbreak{}bin}).

\section{Backend Technologies}

\subsection{Next.js Route Handlers}
\textbf{Purpose:} Serverless API endpoints at \texttt{app/\allowbreak{}api/\allowbreak{}*/\allowbreak{}route.\allowbreak{}ts}.

\subsection{ThingSpeak REST API}
\textbf{Purpose:} IoT cloud data platform; stores and serves field-mapped sensor readings.

\section{IoT Technologies}

\subsection{ESP32 DevKit}
\textbf{Purpose:} Main microcontroller; reads sensors and uploads data to ThingSpeak via Wi-Fi.\\
\textbf{Framework:} Arduino (C++ sketch: \texttt{IOTcode/\allowbreak{}IOTcode.\allowbreak{}ino}).

\subsection{Adafruit BMP085 Library}
\textbf{Purpose:} Arduino driver for the BMP180 pressure and temperature sensor.

\subsection{DHT Sensor Library}
\textbf{Purpose:} Arduino driver for the DHT11 temperature and humidity sensor.

\section{Development Tools}

\begin{table}[H]
\centering
\caption{Development Tools and Versions}
\begin{tabularx}{\linewidth}{|l|l|X|}
\hline
\rowcolor{FloodLight}
\textbf{Tool} & \textbf{Version} & \textbf{Purpose} \\
\hline
Node.js & v20.x & JavaScript runtime \\
npm & v10.x & Package manager \\
TypeScript & 5.x & Static type checking \\
ESLint & 9.x & Code quality linting \\
PostCSS & 4.x & Tailwind CSS processing \\
Google Colab & --- & GPU-accelerated ML training environment \\
Arduino IDE & --- & ESP32 firmware development and flashing \\
\hline
\end{tabularx}
\end{table}

\section{Hosting and Infrastructure}

\begin{table}[H]
\centering
\caption{Hosting and Infrastructure}
\begin{tabularx}{\linewidth}{|l|X|}
\hline
\rowcolor{FloodLight}
\textbf{Service} & \textbf{Role} \\
\hline
Vercel & Next.js application hosting; serverless functions; CDN; CI/CD pipeline \\
ThingSpeak & IoT cloud data storage and REST API \\
MapTiler & Vector tile service for MapLibre GL interactive map \\
\hline
\end{tabularx}
\end{table}
